\section{Implementation}

The implementation is organized as a Python package under \texttt{src/tacc/}. The main modules are:
\begin{itemize}
    \item \texttt{data.py}: loads movement traces and creates deterministic trace-compatible mobility data when raw traces are unavailable;
    \item \texttt{graph.py}: builds handoff graphs, computes graph features, and samples perturbations;
    \item \texttt{demand.py}: generates Zipf-locality demand;
    \item \texttt{evaluate.py}: computes access cost, hit ratios, redundancy, relocation, and objective;
    \item \texttt{placement.py}: implements random, local popularity, global popularity, diversified popularity, and topology-aware greedy placements;
    \item \texttt{rl.py}: implements DQN training and accepted-action refinement;
    \item \texttt{experiment.py}: runs the full experiment and writes CSV and metadata outputs;
    \item \texttt{reporting.py}: generates LaTeX result assets.
\end{itemize}

The experiment is executed with:
\[
\texttt{python3 scripts/run\_experiments.py --output-dir results/final}.
\]
The final run writes \texttt{summary\_metrics.csv}, \texttt{sample\_metrics.csv}, and \texttt{run\_metadata.json}. This separation is important: the summary file supports paper tables, while the sample-level file preserves repeated perturbation measurements for future statistical analysis.

The implementation uses NumPy, pandas, NetworkX, and PyTorch. The LaTeX source is structured as a normal multi-file IEEE-style paper and can be compiled from \texttt{Paper/final\_paper/} on Overleaf or any TeX distribution with IEEEtran support.

\subsection{Reproducibility Checklist}

The experiment is reproducible through the following recorded choices:
\begin{itemize}
    \item random seed: 11;
    \item graph size: 20 AP nodes and 190 weighted edges;
    \item catalog size: 48 objects;
    \item cache capacity: 5 objects per node;
    \item demand model: Zipf with \(\alpha=0.85\) and locality multiplier 2.8;
    \item perturbation rates: 0.00, 0.05, 0.10, and 0.15;
    \item DQN training: 85 episodes and 16 steps per episode;
    \item evaluation: 12 perturbed graph samples per rate.
\end{itemize}

The code writes both summary and sample-level CSV files. This is important because the paper can report mean behavior while future work can compute confidence intervals or run statistical tests from sample-level outputs.
